\section{Recording attribution and its evidence (RQ4)}
\label{sec:provenance}

Sections~\ref{sec:validity} and~\ref{sec:correction} measure how far a detector score is
from an authorship fact. That distance is a property of the score, the genre and the
threshold, and it is known once measured. The problem is that nothing in current
practice carries it: a prevalence figure, a dashboard flag or a moderation decision
records that an artifact was classified as AI-written and discards the grounds. This
section describes a small vocabulary that keeps the two together, and shows the
questions that become answerable once it does.

\subsection{The model}
\label{sec:provenance:model}

\texttt{aiprov} extends PROV-O \citep{lebo2013provo} with one commitment: an authorship
claim is a thing in its own right, carrying its own evidence, rather than a label on an
artifact. An \texttt{aiprov:AuthorshipClaim} points at an \texttt{aiprov:Artifact} and,
where one is named, at an \texttt{aiprov:CodingAgent} identified by tool, model and
version, because the same tool across two model versions is not the same author. Each
claim carries at least one \texttt{aiprov:Evidence}.

Evidence divides by how much it rests on inference. Table~\ref{tab:evidence} lists the
classes.

\begin{table}[t]
\caption{Evidence classes in \texttt{aiprov}, ordered by how far the attribution is from
a statement by the party that made the artifact.}
\label{tab:evidence}
\small
\begin{tabular}{@{}p{0.26\textwidth}p{0.30\textwidth}p{0.34\textwidth}@{}}
\toprule
Class & Where it comes from & What it carries \\
\midrule
Agent account & The platform opened the contribution under its own account &
Agent identity; no error rate \\
\addlinespace
Trailer & A \texttt{Co-authored-by} line naming an agent &
Agent identity; no error rate \\
\addlinespace
Platform metadata & Other fields the platform records &
Agent identity; no error rate \\
\addlinespace
Self-admission & The artifact text names the tool that wrote it &
Agent identity; reliable where present, silent otherwise \\
\addlinespace
Detector & An instrument scored the text above a threshold &
Score, threshold, detector version and weights, and the stratum's false positive rate,
sensitivity and Youden's $J$ \\
\bottomrule
\end{tabular}
\end{table}

The first four rows are declared: authorship is stated rather than inferred, and there
is no error rate to record because nothing was estimated. The last row is the inferred
case, and it is the only one that carries the measurements of
Sections~\ref{sec:validity:rq1} and~\ref{sec:correction}. Attaching them to the evidence
rather than to the detector is what lets a consumer weigh one claim against another: the
same score means different things on a pull request body and on a diff, and the stratum
is what says so.

Two modelling choices are worth stating. The first is that an artifact with no claim is
not an artifact claimed to be human-written. Absence of a claim represents absence of
evidence, and conflating the two is how a detector's silence becomes a positive
assertion about a contributor. The second is that competing claims about one artifact
coexist. A pull request declared by an agent account and flagged by three detectors
carries four claims, and a consumer can ask which of them would survive on its own.

\subsection{Demonstration}
\label{sec:provenance:demo}

The corpus, its detector scores and the stratum error rates of
Tables~\ref{tab:rq1} and~\ref{tab:rq3} are loaded into a graph over a sample of the
corpus. Five competency questions are answered in SPARQL, and each of them turns on the
relation between an attribution and its grounds, so none is answerable from a table of
scores. Two are worth reproducing here.

The first asks which artifacts are attributed on detector evidence alone, in a stratum
whose measured false positive rate exceeds five per cent, with no declared evidence to
corroborate them. These are the attributions a prevalence estimate rests on, and the
query returns them ranked by genre and detector. The second asks for claims whose
supporting evidence comes from a stratum where $J$ is at or below zero. It returns the
claims that carry no information whatever their score, which is a statement no score
threshold can express and which the model makes queryable.

A third question runs in the opposite direction: artifacts an agent is declared to have
written that no detector flagged. That is the recall gap of the detector stack measured
against ground truth, and in the sample graph it is largest on diffs and smallest on
pull request bodies, which matches Table~\ref{tab:rq1}.

The vocabulary, the graph and the queries are in the replication package. The
contribution is the design and its demonstration rather than the graph itself, which is
built over a sample and serves no purpose as a dataset. We do not claim adoption: the
model is worth its length only if it makes the error rates travel with the claims, and
whether that happens depends on the platforms that issue them.
